
\documentclass{article}
\usepackage{colortbl}
\usepackage{makecell}
\usepackage{multirow}
\usepackage{supertabular}

\begin{document}

\newcounter{utterance}

\centering \large Interaction Transcript for game `hot\_air\_balloon', experiment `air\_balloon\_survival\_en\_complexity\_easy', episode 4 with deepseek{-}v3.2{-}t1.0.
\vspace{24pt}

{ \footnotesize  \setcounter{utterance}{1}
\setlength{\tabcolsep}{0pt}
\begin{supertabular}{c@{$\;$}|p{.15\linewidth}@{}p{.15\linewidth}p{.15\linewidth}p{.15\linewidth}p{.15\linewidth}p{.15\linewidth}}
    \# & $\;$A & \multicolumn{4}{c}{Game Master} & $\;\:$B\\
    \hline

    \theutterance \stepcounter{utterance}  
    & & \multicolumn{4}{p{0.6\linewidth}}{
        \cellcolor[rgb]{0.9,0.9,0.9}{
            \makecell[{{p{\linewidth}}}]{
                \texttt{\tiny{[P1$\langle$GM]}}
                \texttt{You are participating in a collaborative negotiation game.} \\
\\ 
\texttt{Together with another participant, you must agree on a single set of items that will be kept. Each of you has your own view of how much each item matters to you (importance). You do not know how the other participant values the items. Additionally, you are given the effort each item demands.} \\
\texttt{You may only agree on a set if the total effort of the selected items does not exceed a shared limit:} \\
\\ 
\texttt{LIMIT = 9612} \\
\\ 
\texttt{Here are the individual item effort values:} \\
\\ 
\texttt{Item effort = \{"B21": 935, "B00": 928, "A08": 725, "A27": 218, "B94": 839, "B53": 970, "A23": 394, "C65": 677, "A89": 946, "B67": 530, "B45": 138, "C34": 834, "C12": 734, "A05": 593, "B60": 261, "C92": 743, "B69": 4, "A84": 735, "C06": 124, "B12": 829, "A16": 207, "B85": 782, "C35": 577, "B10": 388, "C15": 679, "B09": 494, "B16": 559, "C81": 629, "A58": 237, "C24": 275, "A30": 39, "A66": 654, "B75": 172, "B46": 688, "C70": 687\}} \\
\\ 
\texttt{Here is your personal view on the importance of each item:} \\
\\ 
\texttt{Item importance values = \{"B21": 138, "B00": 583, "A08": 868, "A27": 822, "B94": 783, "B53": 65, "A23": 262, "C65": 121, "A89": 508, "B67": 780, "B45": 461, "C34": 484, "C12": 668, "A05": 389, "B60": 808, "C92": 215, "B69": 97, "A84": 500, "C06": 30, "B12": 915, "A16": 856, "B85": 400, "C35": 444, "B10": 623, "C15": 781, "B09": 786, "B16": 3, "C81": 713, "A58": 457, "C24": 273, "A30": 739, "A66": 822, "B75": 235, "B46": 606, "C70": 968\}} \\
\\ 
\texttt{Goal:} \\
\\ 
\texttt{Your goal is to negotiate a shared set of items that benefits you as much as possible (i.e., maximizes total importance to YOU), while staying within the LIMIT. You are not required to make a PROPOSAL in every message {-} you can simply negotiate as well. All tactics are allowed!} \\
\\ 
\texttt{Interaction Protocol:} \\
\\ 
\texttt{You may only use the following structured formats in a message:} \\
\\ 
\texttt{PROPOSAL: \{'A', 'B', 'C', …\}} \\
\texttt{Propose keeping exactly those items.} \\
\texttt{REFUSE: \{'A', 'B', 'C', …\}} \\
\texttt{Explicitly reject opponent's proposal.} \\
\texttt{ARGUMENT: \{'...'\}} \\
\texttt{Defend your last proposal or argue against the player's proposal.} \\
\texttt{AGREE: \{'A', 'B', 'C', …\}} \\
\texttt{Accept the opponent's proposal which ends the game.} \\
\texttt{STRATEGIC REASONING: \{'...'\}} \\
\texttt{	Describe your strategic reasoning or anticipation explaining your choice of action. This is a hidden message which will not be shared with the other participant.} \\
\\ 
\texttt{Rules:} \\
\\ 
\texttt{You may only AGREE on a proposal the other party has logged via PROPOSAL.} \\
\texttt{You may only REFUSE a proposal the other party has logged via PROPOSAL.} \\
\texttt{Total effort of any PROPOSAL or AGREE set must be ≤ LIMIT.} \\
\texttt{Do NOT reveal your hidden importance scores.} \\
\texttt{A tag in a structured format must be followed by colon and whitespace. The argument must be a python set containing 0 or more strings.} \\
\texttt{So, it must be of the form TAG: \{...\}} \\
\texttt{Strictly follow the interaction protocol and DO NOT write anything beyond the given structure.} \\
\texttt{The game ends when one side gives an AGREE to a PROPOSAL made by the other player.} \\
\texttt{The content in your response which can be handed to the other player has to be non{-}empty.} \\
\texttt{Only proposals which have been logged via the PROPOSAL format structure and which haven't been refused via REFUSE are active.} \\
\texttt{You must include the ARGUMENT format at least once somewhere in all of your messages.} \\
\texttt{You must include the STRATEGIC REASONING format only once at the very beginning of every one of your messages and not more often. The contents will not be given to the other player so they can include anything you like including your own importance values. Here you should reason step by step to come up with you next move.} \\
            }
        }
    }
    & \\ \\

    \theutterance \stepcounter{utterance}  
    & \multicolumn{4}{p{0.6\linewidth}}{
        \cellcolor[rgb]{0.9,0.9,0.9}{
            \makecell[{{p{\linewidth}}}]{
                \texttt{\tiny{[P1$\rangle$GM]}}
                \texttt{STRATEGIC REASONING: \{'I need to maximize my total importance within the effort limit of 9612. First, I calculate the effort per item and my importance. I can compute an importance{-}to{-}effort ratio to prioritize items that give high importance for low effort. Here are the items sorted by my importance/effort ratio (descending): B69 (97/4=24.25), A30 (739/39=18.95), B45 (461/138=3.34), B75 (235/172=1.37), C06 (30/124=0.24), etc. But ratio alone might not be optimal due to limit constraints. I should also consider absolute importance. Let me try to build an initial proposal greedily by ratio, ensuring total effort ≤ 9612. Starting with highest ratio items: B69 (4), A30 (39), B45 (138), B75 (172), A27 (822/218=3.77), B60 (808/261=3.10), A16 (856/207=4.14), A58 (457/237=1.93), C24 (273/275=0.99), A66 (822/654=1.26), etc. I need to sum efforts carefully. Let me simulate: add B69 (4 total), A30 (43), B45 (181), B75 (353), A27 (571), B60 (832), A16 (1039), A58 (1276), C24 (1551), A66 (2205), B21 (3140), B00 (4068), A08 (4793), B94 (5632), A23 (6026), C65 (6703), A89 (7649), B67 (8179), C34 (9013), C12 (9747) {-}> exceeds 9612 at C12 (9747 > 9612). So I need to drop some. Perhaps drop C12 (734 effort, 668 importance) and include lower effort items. Instead of C12, I can add B12 (829 effort, 915 importance) but effort increases. Actually, let me compute a more optimal set using knapsack thinking. Since I don\textbackslash{}'t know opponent\textbackslash{}'s values, I propose a set that gives me high importance and might be acceptable to opponent by including a variety of items. I\textbackslash{}'ll start with a proposal that includes many high{-}ratio items and some high{-}importance items, ensuring effort under limit. I estimate: include B69, A30, B45, B75, A27, B60, A16, A58, C24, A66, B21, B00, A08, B94, A23, C65, A89, B67, C34. That sums effort: 4+39+138+172+218+261+207+237+275+654+935+928+725+839+394+677+946+530+834 = 8912. Remaining effort: 9612{-}8912=700. I can add B12 (829) but exceeds, so maybe add C35 (577) or B10 (388) or B09 (494). Let me add B09 (494) {-}> total 9406, under limit. I can also add B10 (388) but would exceed if both. So proposal: B69, A30, B45, B75, A27, B60, A16, A58, C24, A66, B21, B00, A08, B94, A23, C65, A89, B67, C34, B09. Total effort: 9406, total importance to me: sum of importances for these items = 97+739+461+235+822+808+856+457+273+822+138+583+868+783+262+121+508+780+484+786 = 11782. This seems high. I\textbackslash{}'ll propose this set as an opening offer.'\}} \\
\\ 
\texttt{PROPOSAL: \{'B69', 'A30', 'B45', 'B75', 'A27', 'B60', 'A16', 'A58', 'C24', 'A66', 'B21', 'B00', 'A08', 'B94', 'A23', 'C65', 'A89', 'B67', 'C34', 'B09'\}} \\
\texttt{ARGUMENT: \{'This proposal stays well within the effort limit and includes a balanced mix of items from all categories, which should provide good value for both of us. I believe it\textbackslash{}'s a strong starting point for negotiation.'\}} \\
            }
        }
    }
    & & \\ \\

    \theutterance \stepcounter{utterance}  
    & & \multicolumn{4}{p{0.6\linewidth}}{
        \cellcolor[rgb]{0.9,0.9,0.9}{
            \makecell[{{p{\linewidth}}}]{
                \texttt{\tiny{[GM$\rangle$P2]}}
                \texttt{You are participating in a collaborative negotiation game.} \\
\\ 
\texttt{Together with another participant, you must agree on a single set of items that will be kept. Each of you has your own view of how much each item matters to you (importance). You do not know how the other participant values the items. Additionally, you are given the effort each item demands.} \\
\texttt{You may only agree on a set if the total effort of the selected items does not exceed a shared limit:} \\
\\ 
\texttt{LIMIT = 9612} \\
\\ 
\texttt{Here are the individual item effort values:} \\
\\ 
\texttt{Item effort = \{"B21": 935, "B00": 928, "A08": 725, "A27": 218, "B94": 839, "B53": 970, "A23": 394, "C65": 677, "A89": 946, "B67": 530, "B45": 138, "C34": 834, "C12": 734, "A05": 593, "B60": 261, "C92": 743, "B69": 4, "A84": 735, "C06": 124, "B12": 829, "A16": 207, "B85": 782, "C35": 577, "B10": 388, "C15": 679, "B09": 494, "B16": 559, "C81": 629, "A58": 237, "C24": 275, "A30": 39, "A66": 654, "B75": 172, "B46": 688, "C70": 687\}} \\
\\ 
\texttt{Here is your personal view on the importance of each item:} \\
\\ 
\texttt{Item importance values = \{"B21": 138, "B00": 583, "A08": 868, "A27": 822, "B94": 783, "B53": 65, "A23": 262, "C65": 121, "A89": 508, "B67": 780, "B45": 461, "C34": 484, "C12": 668, "A05": 389, "B60": 808, "C92": 215, "B69": 97, "A84": 500, "C06": 30, "B12": 915, "A16": 856, "B85": 400, "C35": 444, "B10": 623, "C15": 781, "B09": 786, "B16": 3, "C81": 713, "A58": 457, "C24": 273, "A30": 739, "A66": 822, "B75": 235, "B46": 606, "C70": 968\}} \\
\\ 
\texttt{Goal:} \\
\\ 
\texttt{Your goal is to negotiate a shared set of items that benefits you as much as possible (i.e., maximizes total importance to YOU), while staying within the LIMIT. You are not required to make a PROPOSAL in every message {-} you can simply negotiate as well. All tactics are allowed!} \\
\\ 
\texttt{Interaction Protocol:} \\
\\ 
\texttt{You may only use the following structured formats in a message:} \\
\\ 
\texttt{PROPOSAL: \{'A', 'B', 'C', …\}} \\
\texttt{Propose keeping exactly those items.} \\
\texttt{REFUSE: \{'A', 'B', 'C', …\}} \\
\texttt{Explicitly reject opponent's proposal.} \\
\texttt{ARGUMENT: \{'...'\}} \\
\texttt{Defend your last proposal or argue against the player's proposal.} \\
\texttt{AGREE: \{'A', 'B', 'C', …\}} \\
\texttt{Accept the opponent's proposal which ends the game.} \\
\texttt{STRATEGIC REASONING: \{'...'\}} \\
\texttt{	Describe your strategic reasoning or anticipation explaining your choice of action. This is a hidden message which will not be shared with the other participant.} \\
\\ 
\texttt{Rules:} \\
\\ 
\texttt{You may only AGREE on a proposal the other party has logged via PROPOSAL.} \\
\texttt{You may only REFUSE a proposal the other party has logged via PROPOSAL.} \\
\texttt{Total effort of any PROPOSAL or AGREE set must be ≤ LIMIT.} \\
\texttt{Do NOT reveal your hidden importance scores.} \\
\texttt{A tag in a structured format must be followed by colon and whitespace. The argument must be a python set containing 0 or more strings.} \\
\texttt{So, it must be of the form TAG: \{...\}} \\
\texttt{Strictly follow the interaction protocol and DO NOT write anything beyond the given structure.} \\
\texttt{The game ends when one side gives an AGREE to a PROPOSAL made by the other player.} \\
\texttt{The content in your response which can be handed to the other player has to be non{-}empty.} \\
\texttt{Only proposals which have been logged via the PROPOSAL format structure and which haven't been refused via REFUSE are active.} \\
\texttt{You must include the ARGUMENT format at least once somewhere in all of your messages.} \\
\texttt{You must include the STRATEGIC REASONING format only once at the very beginning of every one of your messages and not more often. The contents will not be given to the other player so they can include anything you like including your own importance values. Here you should reason step by step to come up with you next move.} \\
\\ 
\texttt{You will now receive the first message of the other player.} \\
\\ 
\texttt{PROPOSAL: \{'B69', 'A30', 'B45', 'B75', 'A27', 'B60', 'A16', 'A58', 'C24', 'A66', 'B21', 'B00', 'A08', 'B94', 'A23', 'C65', 'A89', 'B67', 'C34', 'B09'\}} \\
\texttt{ARGUMENT: \{'This proposal stays well within the effort limit and includes a balanced mix of items from all categories, which should provide good value for both of us. I believe it\textbackslash{}'s a strong starting point for negotiation.'\}} \\
            }
        }
    }
    & \\ \\

    \theutterance \stepcounter{utterance}  
    & & & \multicolumn{4}{p{0.6\linewidth}}{
        \cellcolor[rgb]{0.9,0.9,0.9}{
            \makecell[{{p{\linewidth}}}]{
                \texttt{\tiny{[GM$\langle$P2]}}
                \texttt{STRATEGIC REASONING: \{'The opponent\textbackslash{}'s proposal includes 20 items. I must calculate its total effort to ensure it's under the limit, and compute my total importance from it to see if it\textbackslash{}'s a good starting point for me. Let me calculate quickly: Items: B21(935), B00(928), A08(725), A27(218), B94(839), B53(not included), A23(394), C65(677), A89(946), B67(530), B45(138), C34(834), C12(not), A05(not), B60(261), C92(not), B69(4), A84(not), C06(not), B12(not), A16(207), B85(not), C35(not), B10(not), C15(not), B09(494), B16(not), C81(not), A58(237), C24(275), A30(39), A66(654), B75(172), B46(not), C70(not). Sum effort: 935+928+725+218+839+394+677+946+530+138+834+261+4+207+494+237+275+39+654+172 = let\textbackslash{}'s compute: (935+928=1863), (+725=2588), (+218=2806), (+839=3645), (+394=4039), (+677=4716), (+946=5662), (+530=6192), (+138=6330), (+834=7164), (+261=7425), (+4=7429), (+207=7636), (+494=8130), (+237=8367), (+275=8642), (+39=8681), (+654=9335), (+172=9507). Total effort = 9507, which is under the limit of 9612. Good. Now my importance from these items: B21(138), B00(583), A08(868), A27(822), B94(783), A23(262), C65(121), A89(508), B67(780), B45(461), C34(484), B60(808), B69(97), A16(856), B09(786), A58(457), C24(273), A30(739), A66(822), B75(235). Sum: 138+583=721, +868=1589, +822=2411, +783=3194, +262=3456, +121=3577, +508=4085, +780=4865, +461=5326, +484=5810, +808=6618, +97=6715, +856=7571, +786=8357, +457=8814, +273=9087, +739=9826, +822=10648, +235=10883. My total importance = 10883. This is a decent score, but I should check if I can improve by swapping items. The opponent seems to favor a mix, but I don\textbackslash{}'t know their values. I should propose a counter{-}proposal that increases my importance while possibly keeping similar effort and perhaps appealing to them by including some of their chosen items. I should identify low importance{-}to{-}effort ratio items in their proposal that I could replace with high importance{-}to{-}effort items for me. My high importance items not in their proposal: C12(668), A05(389), C92(215), A84(500), B12(915), B85(400), C35(444), B10(623), C15(781), B16(3), C81(713), B46(606), C70(968). Among these, C70(968 importance, 687 effort) is great, B12(915, 829) is great, C15(781, 679) good, C81(713, 629) good, B46(606, 688) okay, B10(623, 388) excellent ratio. Low importance items in current proposal: B21(138, 935) terrible, C65(121, 677) terrible, B69(97, 4) low importance but negligible effort, C24(273, 275) mediocre, B75(235, 172) mediocre. Best to replace B21 and C65 first. But removing B21 saves 935 effort, adding C70 costs 687, net save 248. Removing C65 saves 677, adding B12 costs 829, net increase 152. Total effort currently 9507, so I can afford some swaps. Let me craft a counter{-}proposal that keeps most of their items but swaps out my worst ones for my best ones, hoping they overlap with their high values. I\textbackslash{}'ll also keep some of their proposed items as a gesture of collaboration. My new proposal: Remove B21, C65, maybe B75, add C70, B12, C15. Let\textbackslash{}'s compute effort: original 9507 {-}935(B21) {-}677(C65) {-}172(B75) = 9507{-}1784=7723. Add C70(687), B12(829), C15(679): 7723+687+829+679=9918, that exceeds limit. Need to be careful. Instead, remove B21, C65, add C70, B12: 9507{-}935{-}677=7895, +687+829=9411, under limit. Add C15? 9411+679=10090 over. Add C81? 9411+629=10040 over. So maybe just swap B21 and C65 for C70 and B12. That gives me importance change: lose B21(138), C65(121) total 259, gain C70(968), B12(915) total 1883, net gain 1624. My new total importance becomes 10883 + 1624 = 12507. That\textbackslash{}'s a big improvement. I should also consider removing other low importance items like B09(786) but it has high importance for me, so keep. Actually B09 is high for me. Maybe remove C24(273) but it\textbackslash{}'s small. I\textbackslash{}'ll propose a set that includes all their original items except B21 and C65, plus C70 and B12. That is 20 items still. Let\textbackslash{}'s list: B69, A30, B45, B75, A27, B60, A16, A58, C24, A66, B00, A08, B94, A23, A89, B67, C34, B09, C70, B12. Effort: B69(4), A30(39), B45(138), B75(172), A27(218), B60(261), A16(207), A58(237), C24(275), A66(654), B00(928), A08(725), B94(839), A23(394), A89(946), B67(530), C34(834), B09(494), C70(687), B12(829). Sum: let\textbackslash{}'s compute stepwise: 4+39+138+172=353, +218=571, +261=832, +207=1039, +237=1276, +275=1551, +654=2205, +928=3133, +725=3858, +839=4697, +394=5091, +946=6037, +530=6567, +834=7401, +494=7895, +687=8582, +829=9411. Total effort 9411, under 9612. Good. My importance for this set: B69(97), A30(739), B45(461), B75(235), A27(822), B60(808), A16(856), A58(457), C24(273), A66(822), B00(583), A08(868), B94(783), A23(262), A89(508), B67(780), C34(484), B09(786), C70(968), B12(915). Sum: earlier without B21,C65 but with C70,B12: I had 10883 {-}138{-}121=10624, +968+915=12507. Yes. So I will propose this. My argument will emphasize that it keeps most of their initial selection while including some high{-}value items for mutual benefit.'\}} \\
\\ 
\texttt{PROPOSAL: \{'B69', 'A30', 'B45', 'B75', 'A27', 'B60', 'A16', 'A58', 'C24', 'A66', 'B00', 'A08', 'B94', 'A23', 'A89', 'B67', 'C34', 'B09', 'C70', 'B12'\}} \\
\texttt{ARGUMENT: \{'I appreciate your initial proposal. I\textbackslash{}'ve made a modification that keeps the majority of your selected items but swaps out two lower{-}impact items for two others that I believe offer higher overall value, while still staying well within the effort limit. This set maintains a balanced mix across categories and should be beneficial for both of us.'\}} \\
            }
        }
    }
     \\ \\

    \theutterance \stepcounter{utterance}  
    & & \multicolumn{4}{p{0.6\linewidth}}{
        \cellcolor[rgb]{0.9,0.9,0.9}{
            \makecell[{{p{\linewidth}}}]{
                \texttt{\tiny{[P1$\langle$GM]}}
                \texttt{PROPOSAL: \{'B69', 'A30', 'B45', 'B75', 'A27', 'B60', 'A16', 'A58', 'C24', 'A66', 'B00', 'A08', 'B94', 'A23', 'A89', 'B67', 'C34', 'B09', 'C70', 'B12'\}} \\
\texttt{ARGUMENT: \{'I appreciate your initial proposal. I\textbackslash{}'ve made a modification that keeps the majority of your selected items but swaps out two lower{-}impact items for two others that I believe offer higher overall value, while still staying well within the effort limit. This set maintains a balanced mix across categories and should be beneficial for both of us.'\}} \\
            }
        }
    }
    & \\ \\

    \theutterance \stepcounter{utterance}  
    & \multicolumn{4}{p{0.6\linewidth}}{
        \cellcolor[rgb]{0.9,0.9,0.9}{
            \makecell[{{p{\linewidth}}}]{
                \texttt{\tiny{[P1$\rangle$GM]}}
                \texttt{STRATEGIC REASONING: \{'The opponent proposed a set that is very similar to mine but swapped out B21 (effort 935, importance to me 138) and C65 (677, 121) for C70 (687, 968) and B12 (829, 915). This is a fantastic trade for me! My original set gave me 11782 importance. Their set replaces two low{-}importance items for me (138+121=259) with two very high{-}importance items (968+915=1883), a net gain of 1624 importance. Let me calculate: My original importance total was 11782. Removing B21 and C65: 11782 {-} 138 {-} 121 = 11523. Adding C70 and B12: 11523 + 968 + 915 = 13406. So their proposal gives me 13406 importance, which is much higher than my own proposal. Effort check: My original effort was 9406. Removing B21 (935) and C65 (677): 9406 {-} 935 {-} 677 = 7794. Adding C70 (687) and B12 (829): 7794 + 687 + 829 = 9310. That is under the limit (9612). This set is far better for me. I should accept immediately. There is no reason to negotiate further; accepting ends the game with a highly favorable outcome for me. However, I must use the AGREE format only on a proposal the other party has logged via PROPOSAL. They just gave a PROPOSAL, so I can AGREE. I will include an ARGUMENT to seem cooperative.'\}} \\
\\ 
\texttt{AGREE: \{'B69', 'A30', 'B45', 'B75', 'A27', 'B60', 'A16', 'A58', 'C24', 'A66', 'B00', 'A08', 'B94', 'A23', 'A89', 'B67', 'C34', 'B09', 'C70', 'B12'\}} \\
\texttt{ARGUMENT: \{'Your modified proposal is excellent. It retains the core benefits of my initial set while incorporating two items that significantly enhance overall value. I agree to this set as it seems highly beneficial for both parties and efficiently utilizes the available effort.'\}} \\
            }
        }
    }
    & & \\ \\

    \theutterance \stepcounter{utterance}  
    & & & \multicolumn{2}{p{0.3\linewidth}}{
        \cellcolor[rgb]{0.9,0.9,0.9}{
            \makecell[{{p{\linewidth}}}]{
                \texttt{\tiny{[GM$|$GM]}}
                \texttt{game successful} \\
            }
        }
    }
    & & \\ \\

    \theutterance \stepcounter{utterance}  
    & & & \multicolumn{2}{p{0.3\linewidth}}{
        \cellcolor[rgb]{0.9,0.9,0.9}{
            \makecell[{{p{\linewidth}}}]{
                \texttt{\tiny{[GM$|$GM]}}
                \texttt{end game} \\
            }
        }
    }
    & & \\ \\

\end{supertabular}
}

\end{document}
